%% kapa: exam=1 round=21 year=2010 period=1 form=A subject=economics subject_ko=경제원론 seq=29 number=69
%% source: Q-Net 감정평가사 기출문제 (cst00309) artlSeq=5024181 file=extracted/1008/5024181/2010년 제21회 감정평가사 1차시험 기출문제/1교시.hwp
%% conversion: kordoc 4.24.51 -> tools/gen1.py (auto); review=proofread-sonnet; flags=tikz
\begin{question}{69}
다음 그림에서처럼 X재 시장에서의 균형은 점 E에서 이루어지고 있다. 정부가 단위당 1,000원의 세금을 부과한 결과 X재의 세후 가격은 5,400원으로 상승하였다. 세금 부과로 인한 경제적 순손실(deadweight loss)은 얼마인가?
\begin{center}
\begin{tikzpicture}[font=\small, x=1cm, y=1cm]
  \draw (0,0) -- (0,4.7);
  \draw (0,0) -- (5.3,0);
  \node[above] at (0.1,4.7) {\textbf{가격}};
  \node[below] at (5.2,-0.1) {\textbf{수량}};
  \draw (0,3.92) -- (4.75,1.2);
  \draw (0,0.70) -- (4.42,3.52);
  \draw[dashed] (0,3.11) node[left] {5,400} -- (1.40,3.11) -- (1.40,0) node[below] {4,000};
  \draw[dashed] (0,2.40) node[left] {5,000} -- (2.65,2.40) -- (2.65,0) node[below] {5,000};
  \draw[dashed] (0,1.59) node[left] {4,400} -- (1.40,1.59);
  \fill (2.65,2.40) circle (2pt);
  \node at (2.67,2.78) {\textit{E}};
  \node[right] at (4.5,3.55) {공급};
  \node[right] at (4.8,1.15) {수요};
\end{tikzpicture}
\end{center}
\begin{choices}
  \item 400,000원
  \item 500,000원
  \item 600,000원
  \item 700,000원
  \item 1,000,000원
\end{choices}
\end{question}
\answer{2}
